\subsection{Astronomical Alert Brokers}
\label{sec:related:brokers}

The ZTF alert stream has prompted the development of several 
automated classification brokers. ALeRCE \citep{Forster2021} 
operates a two-stage pipeline: a stamp classifier for early 
real/bogus filtering and a light curve classifier 
\citep{Sanchez-Saez2021} for multi-class transient 
classification using Random Forest on extracted photometric 
features. Fink \citep{Moller2021} provides a modular 
science portal with multiple classifiers operating 
independently, including the SuperNNova-based binary 
classifier \citep{Moller2019} and a Random Forest 
SN Ia discriminator. ANTARES \citep{Narayan2018} focuses 
on real-time annotation and filtering. Lasair 
\citep{Smith2019} provides alert stream access with 
community-contributed classifiers.

All major brokers provide classification probabilities 
as outputs. None, to our knowledge, publish calibration 
metrics such as ECE or reliability diagrams for their 
production classifiers. This gap motivates the present work.

\subsection{Calibration in Machine Learning}
\label{sec:related:calibration}

The calibration of neural network classifiers has received 
substantial attention since \citet{Guo2017} demonstrated 
that modern deep networks are systematically overconfident 
and that temperature scaling provides an effective post-hoc 
remedy. Subsequent work has examined calibration under 
distribution shift \citep{Ovadia2019}, for binary 
classifiers \citep{Kull2017}, and in the context of 
selective prediction \citep{Geifman2017}.

Calibration of Random Forest classifiers is less studied 
but well understood at a mechanistic level. 
\citet{Niculescu-Mizil2005} demonstrated that RF produces 
systematically distorted probability estimates — 
underconfident for well-discriminated classes and poorly 
calibrated near class boundaries — and that Platt scaling 
and isotonic regression provide effective corrections. 
\citet{Zadrozny2002} introduced isotonic regression as 
a nonparametric post-hoc calibration method applicable 
to any classifier.

In astronomy, calibration has been studied in the context 
of photometric redshift estimation \citep{Bordoloi2010, 
Schmidt2020}, where reliable uncertainty quantification 
directly affects cosmological inference. Calibration of 
transient classifiers has not, to our knowledge, been 
systematically evaluated on real spectroscopic samples 
prior to this work.

\subsection{Uncertainty Quantification for Transient Classification}
\label{sec:related:uq}

Uncertainty quantification for astronomical transient 
classification has been explored primarily through 
Bayesian and ensemble approaches. \citet{Möller2020} 
applied Bayesian neural networks to supernova 
photometric classification, obtaining uncertainty 
estimates alongside class predictions. \citet{Villar2019} 
used deep ensembles for the PLAsTiCC challenge, 
demonstrating that ensemble disagreement correlates 
with classification difficulty.

These works focus on model uncertainty (epistemic 
uncertainty from limited training data) rather than 
calibration (the alignment between predicted confidence 
and empirical accuracy). The two concepts are related 
but distinct: a well-calibrated classifier need not 
have low uncertainty, and a classifier with well-quantified 
uncertainty need not be calibrated. Our focus on ECE 
and reliability diagrams addresses the calibration 
question directly, complementing prior uncertainty 
quantification work.

\subsection{Learning to Defer}
\label{sec:related:l2d}

Learning-to-defer \citep[L2D;][]{Madras2018} is a 
framework for training joint classifier-deferral systems 
that learn when to route decisions to a human expert. 
\citet{Mozannar2020} extended L2D to handle consistent 
surrogate losses, and \citet{Verma2022} developed 
calibrated L2D systems that provide reliable confidence 
estimates alongside deferral decisions.

L2D has been applied in medical imaging 
\citep{Tschandl2019}, content moderation, and 
autonomous driving, where human-AI collaboration 
under resource constraints is operationally critical. 
To our knowledge, L2D has not been applied to 
astronomical transient classification. The present 
calibration study establishes the empirical foundation 
for this application: demonstrating that production 
classifiers are miscalibrated motivates the need for 
explicit deferral mechanisms, and providing post-hoc 
calibration methods identifies the broker most 
suitable for downstream L2D integration.